\chapter{Kết quả Thực nghiệm và Phân tích Kỹ thuật}
\label{chap:results}

Chương này trình bày kết quả thực nghiệm thu thập trên cụm \textbf{MongoDB 7.0} triển khai tại Microsoft Azure theo đúng ba thực nghiệm \texttt{TN1}\textendash \texttt{TN3} của Chương~\ref{chap:methodology}. Mỗi câu hỏi nghiên cứu được gói gọn trong \textbf{đúng một bảng kết quả chính} chỉ chứa các chỉ số lõi, đi kèm biểu đồ trực quan và phân tích kỹ thuật dựa trên dữ liệu tài nguyên hệ thống (\texttt{sar}, \texttt{serverStatus}, \texttt{explain}). Toàn bộ dữ liệu CSV/JSON thô được lưu trữ công khai tại kho GitHub \cite{ProjectRepo}.

\resultplaceholder{Trạng thái chương}{Đây là bộ khung trung lập gồm đúng ba bảng kết quả cho Chương 4 (không điền trước kết quả kỳ vọng). Sau khi hoàn tất chạy $3$ pha trên Azure ($\sim 2$\textendash $2.5$ giờ) và xác nhận CPU máy khách $< 75\%$, thay các ký hiệu \texttt{--} bằng số liệu trung bình qua $3$ lần lặp và chèn biểu đồ tương ứng.}

\section{Kết quả RQ1: Mở rộng Đọc theo Quy mô và Số luồng (\texttt{TN1})}

Bảng~\ref{tab:rq1-results} đối chiếu trực tiếp độ trễ \texttt{p95} và thông lượng \texttt{QPS} của truy vấn \texttt{Q1} giữa máy đơn (\texttt{SA}) và cụm $2$ máy (\texttt{SH-hash}) qua bốn mốc quy mô $100\text{K} \rightarrow 3\text{M}$ và ba mức luồng $1 / 4 / 8$.

\begin{table}[!htb]
\centering
\caption{So sánh hiệu năng đọc \texttt{Q1} giữa \texttt{SA} và \texttt{SH-hash} qua bốn mốc quy mô và ba mức luồng (\texttt{TN1}).}
\label{tab:rq1-results}
\renewcommand{\arraystretch}{1.18}
\begin{tabularx}{\linewidth}{l c C{2.4cm} C{2.6cm} C{2.3cm} C{2.5cm}}
\toprule
\textbf{Quy mô} & \textbf{Số luồng} & \textbf{\texttt{SA} \texttt{p95} (ms)} & \textbf{\texttt{SH-hash} \texttt{p95} (ms)} & \textbf{\texttt{SA} \texttt{QPS}} & \textbf{\texttt{SH-hash} \texttt{QPS}} \\
\midrule
\multirow{3}{*}{\texttt{100K}} & 1 & \textendash{} & \textendash{} & \textendash{} & \textendash{} \\
                               & 4 & \textendash{} & \textendash{} & \textendash{} & \textendash{} \\
                               & 8 & \textendash{} & \textendash{} & \textendash{} & \textendash{} \\
\midrule
\multirow{3}{*}{\texttt{500K}} & 1 & \textendash{} & \textendash{} & \textendash{} & \textendash{} \\
                               & 4 & \textendash{} & \textendash{} & \textendash{} & \textendash{} \\
                               & 8 & \textendash{} & \textendash{} & \textendash{} & \textendash{} \\
\midrule
\multirow{3}{*}{\texttt{1M}}   & 1 & \textendash{} & \textendash{} & \textendash{} & \textendash{} \\
                               & 4 & \textendash{} & \textendash{} & \textendash{} & \textendash{} \\
                               & 8 & \textendash{} & \textendash{} & \textendash{} & \textendash{} \\
\midrule
\multirow{3}{*}{\texttt{3M}}   & 1 & \textendash{} & \textendash{} & \textendash{} & \textendash{} \\
                               & 4 & \textendash{} & \textendash{} & \textendash{} & \textendash{} \\
                               & 8 & \textendash{} & \textendash{} & \textendash{} & \textendash{} \\
\bottomrule
\end{tabularx}
\end{table}

\resultplaceholder{Biểu đồ Hình 4.1 \& Hình 4.2 (RQ1)}{
  \begin{enumerate}
    \item \textbf{Hình 4.1:} Đồ thị đường \texttt{p95} (ms) theo quy mô $100\text{K} \rightarrow 3\text{M}$ so sánh \texttt{SA} với \texttt{SH-hash}, chỉ rõ mốc giao cắt (\textit{break-even point}) kèm trích dẫn số trang đọc đĩa (\texttt{pages\_read\_into\_cache}) và đọc SSD từ \texttt{sar} để giải thích nguyên nhân.
    \item \textbf{Hình 4.2:} Biểu đồ cột so sánh thông lượng \texttt{QPS} ở $1 / 4 / 8$ luồng tại mốc \texttt{3M} giữa \texttt{SA} và \texttt{SH-hash}.
  \end{enumerate}
}

\section{Kết quả RQ2: Độ đều Dữ liệu và Tốc độ Ghi mới (\texttt{TN2})}

Bảng~\ref{tab:rq2-results} tổng hợp tỷ lệ phân bố dung lượng dữ liệu sau khi Balancer ổn định (từ \texttt{getShardDistribution()}), tỷ lệ lượt ghi trực tuyến rơi vào từng shard (từ \texttt{opcounters.insert}), tốc độ ghi (\texttt{inserts/s}) và độ trễ ghi \texttt{p95} của ba cấu hình tại mốc \texttt{1M}.

\begin{table}[!htb]
\centering
\caption{Phân bố dữ liệu và hiệu năng ghi trực tuyến (\texttt{START = now()}, $4$ luồng, $60\text{s}$) tại mốc \texttt{1M} (\texttt{TN2}).}
\label{tab:rq2-results}
\renewcommand{\arraystretch}{1.24}
\resizebox{\linewidth}{!}{%
\begin{tabular}{l c c r r}
\toprule
\textbf{Cấu hình} & \textbf{Dữ liệu: Máy 1 / Máy 2 (\%)} & \textbf{Lượng ghi mới: Máy 1 / Máy 2 (\%)} & \textbf{Tốc độ ghi (\texttt{inserts/s})} & \textbf{Độ trễ ghi \texttt{p95} (ms)} \\
\midrule
\texttt{SA} \textit{(1 máy)}                & $100\% / 0\%$ & $100\% / 0\%$ & \textendash{} & \textendash{} \\
\texttt{SH-hash} \textit{(băm \texttt{PATIENT})}  & \textendash{}\% / \textendash{}\% & \textendash{}\% / \textendash{}\% & \textendash{} & \textendash{} \\
\texttt{SH-range} \textit{(khoảng \texttt{START})} & \textendash{}\% / \textendash{}\% & \textendash{}\% / \textendash{}\% & \textendash{} & \textendash{} \\
\bottomrule
\end{tabular}%
}
\end{table}

\resultplaceholder{Biểu đồ Hình 4.3 (RQ2)}{Chèn biểu đồ so sánh trực quan tỷ lệ phân bổ lượt ghi giữa Máy 1 và Máy 2 cùng tốc độ ghi \texttt{inserts/s} của \texttt{SA}, \texttt{SH-hash} và \texttt{SH-range}. Phân tích hiện tượng điểm nóng ghi (\textit{Write Hotspot}) của \texttt{SH-range}, vai trò của hai ổ SSD độc lập đối với tăng tốc ghi của \texttt{SH-hash} và ngưỡng dừng $96\text{ MB}$ của Balancer MongoDB 7.0.}

\section{Kết quả RQ3: Tỉa mảnh và Chi phí Truy vấn Xuyên mảnh (\texttt{TN3})}

Bảng~\ref{tab:rq3-results} đối chiếu trực tiếp số máy bị hỏi (đọc từ \texttt{explain("executionStats")}), độ trễ \texttt{p95} và lưu lượng mạng thực qua \texttt{mongos} (đã trừ nền $15\text{s}$ nhàn rỗi) của ba truy vấn \texttt{Q1}, \texttt{Q3}, \texttt{Q4} giữa \texttt{SH-hash} và \texttt{SH-range} tại mốc chuẩn \texttt{1M}.

\begin{table}[!htb]
\centering
\caption{Đối chiếu số máy bị hỏi, độ trễ \texttt{p95} và lưu lượng mạng qua \texttt{mongos} tại mốc \texttt{1M} (\texttt{TN3}).}
\label{tab:rq3-results}
\renewcommand{\arraystretch}{1.24}
\resizebox{\linewidth}{!}{%
\begin{tabular}{l c r r c r}
\toprule
\textbf{Truy vấn} & \textbf{\texttt{SH-hash}: Máy hỏi} & \textbf{\texttt{SH-hash}: \texttt{p95} (ms)} & \textbf{\texttt{SH-hash}: Mạng (KB/req)} & \textbf{\texttt{SH-range}: Máy hỏi} & \textbf{\texttt{SH-range}: \texttt{p95} (ms)} \\
\midrule
\texttt{Q1} \textit{(1 bệnh nhân \textemdash  từ \texttt{TN1})} & \textendash{} / 2 & \textendash{} & \textendash{} & \textemdash{} & \textemdash{} \\
\texttt{Q3} \textit{(Xuyên BN, kết quả nhỏ)}                     & \textendash{} / 2 & \textendash{} & \textendash{} & \textendash{} / 2 & \textendash{} \\
\texttt{Q4} \textit{(Xuyên BN, kết quả lớn)}                     & \textendash{} / 2 & \textendash{} & \textendash{} & \textendash{} / 2 & \textendash{} \\
\bottomrule
\end{tabular}%
}
\end{table}

\resultplaceholder{Biểu đồ Hình 4.4 (RQ3 \textemdash  Độ dốc 100K $\rightarrow$ 3M trên \texttt{SH-hash})}{Chèn đồ thị đường biểu diễn tốc độ tăng độ trễ \texttt{p95} của \texttt{Q1} (1 máy), \texttt{Q3} (2 máy, kết quả nhỏ) và \texttt{Q4} (2 máy, kết quả lớn) trên \texttt{SH-hash} qua bốn mốc $100\text{K} \rightarrow 500\text{K} \rightarrow 1\text{M} \rightarrow 3\text{M}$. Kết hợp phân tích \texttt{mergeType} của \texttt{Q3}, sự đánh đổi \textit{Shard Pruning} vs \textit{Parallel Scan} giữa \texttt{SH-hash} và \texttt{SH-range}, và khác biệt băng thông mạng giữa \texttt{Q3} với \texttt{Q4}.}
